\documentclass[11pt]{article}
\usepackage[margin=1in]{geometry}
\usepackage{amsmath,amssymb,amsthm}
\usepackage{graphicx}
\usepackage{hyperref}

\newtheorem{theorem}{Theorem}
\newtheorem{proposition}{Proposition}
\newtheorem{lemma}{Lemma}
\newtheorem{remark}{Remark}

\newcommand{\E}{\mathbb E}
\newcommand{\Pp}{\mathbb P}
\newcommand{\R}{\mathbb R}
\newcommand{\cB}{\mathcal B}
\newcommand{\norm}[1]{\left\lVert #1\right\rVert}
\newcommand{\Lq}{L^q(S)}
\newcommand{\PhiF}{\Phi}

\title{The \(L^q\), \(p=1\) Endpoint via Poissonization}
\author{Agent \texttt{agent\_lane\_19}}
\date{2026-06-16}

\begin{document}
\maketitle

\section*{Status}

This packet gives a candidate partial solution to Remark 12.10 of
Yaroslavtsev's paper \emph{Local characteristics and tangency of
vector-valued martingales} \cite{Yaroslavtsev2019}.  It is partial relative
to the source's broad UMD-space problem, but it closes the full \(p=1\)
endpoint for the concrete class \(L^q(S)\), \(1<q<\infty\).

Earlier packets in this run prove the \(L^q(S)\)-valued case for
\(1<p,q<\infty\) and prove the \(p=1\) endpoint for \(1<q\le2\) by a
Bernstein-function argument.  The new contribution here is the missing
\(q>2\) square-function comparison, proved by a Poissonization equivalence for
positive \(L^{q/2}\)-valued sums.

\section*{Source Problem}

The source is arXiv:1907.11588, PDF page 105, Remark 12.10.  The source states
that the general local-martingale case is open and reduces the missing issue
to an independent-increment comparison.

\begin{center}
\includegraphics[width=\textwidth]{figures/open_problem_crop.png}
\end{center}

The continuation states the aggregate-domination formulation and notes that
even the symmetric case was not known to the author.

\begin{center}

\includegraphics[width=\textwidth]{figures/open_problem_continuation_crop.png}
\end{center}

\section*{Main Results}

For a finite independent family \(x=(x_i)\) of \(L^q(S)\)-valued random
variables, put
\[
  \mathsf{SF}_q(x)
  =
  \E\left\|
      \left(\sum_i |x_i|^2\right)^{1/2}
    \right\|_{L^q(S)}.
\]
Let
\[
   \mu_x=\sum_i \mathcal L(x_i)|_{\Lq\setminus\{0\}}
\]
be the aggregate occurrence measure.

\begin{theorem}[Endpoint square-function comparison]
\label{thm:square}
Let \(1<q<\infty\).  There is a constant \(C_q\) such that for every measure
space \(S\) and every two finite independent families \(x=(x_i)\),
\(y=(y_j)\) of \(L^q(S)\)-valued random variables,
\[
  \mu_x\le \mu_y
  \quad\Longrightarrow\quad
  \mathsf{SF}_q(x)\le C_q\,\mathsf{SF}_q(y).
\]
\end{theorem}

\begin{theorem}[General local \(L^q\), \(p=1\)]
\label{thm:main}
Let \(1<q<\infty\), let \(S\) be a measure space, and put \(X=L^q(S)\).  If
\(N\) and \(M\) are \(X\)-valued local martingales and \(N\) is
characteristically dominated by \(M\) in the sense of
\cite[Definition 12.9]{Yaroslavtsev2019}, then
\[
  \E\sup_{t\ge0}\norm{N_t}_{L^q(S)}
  \le
  C_q
  \E\sup_{t\ge0}\norm{M_t}_{L^q(S)}.
\]
\end{theorem}

\section*{Idea of the Proof}

For \(1<q\le2\), the previous packet proves Theorem \ref{thm:square} by a
positive multivariate Bernstein representation.  For \(q>2\), put
\[
  r=q/2>1,\qquad \alpha=1/2.
\]
The square-function comparison becomes a comparison for independent positive
\(L^r(S)\)-valued variables:
\[
  \E\norm{\sum_i V_i}_{L^r}^{\alpha}
  \lesssim_r
  \E\norm{\sum_j W_j}_{L^r}^{\alpha}
\]
under aggregate occurrence-measure domination.

The key new observation is that this expectation is equivalent to the same
expectation for the compound-Poisson positive sum whose intensity is the
aggregate occurrence measure.  The compound-Poisson sum depends only on that
aggregate measure.  Therefore aggregate domination becomes immediate: if
\(\nu\le\mu\), the Poisson sum with intensity \(\mu\) is the sum of the
Poisson sum with intensity \(\nu\) and an independent positive remainder.

\section*{Thinning Lemmas}

Throughout this section \(r>1\), \(\alpha=1/2\), and
\[
  \PhiF(f)=\norm{f}_{L^r(S)}^\alpha,\qquad f\in L^r(S)_+.
\]
The functional \(\PhiF\) is monotone on the positive cone and subadditive.

\begin{lemma}[Scalar positive Rosenthal]
\label{lem:positive-rosenthal}
For every \(r>1\) there is a constant \(C_r\) such that for all finite
independent nonnegative scalar random variables \(Z_1,\ldots,Z_n\),
\[
  \E\left(\sum_i Z_i\right)^r
  \le
  C_r\left[
      \left(\sum_i \E Z_i\right)^r
      +
      \sum_i \E Z_i^r
    \right].
\]
\end{lemma}

\begin{proof}
For \(1<r\le2\) this follows from the von Bahr--Esseen inequality applied to
\(Z_i-\E Z_i\), plus the triangle inequality in \(L^r\).  For \(r>2\), the
usual scalar Rosenthal inequality gives
\[
  \E\left|\sum_i(Z_i-\E Z_i)\right|^r
  \le
  C_r\left[
    \left(\sum_i \E Z_i^2\right)^{r/2}
    +
    \sum_i \E Z_i^r
  \right].
\]
Let \(A=\sum_i\E Z_i\) and \(B=\sum_i\E Z_i^r\).  By interpolation,
\[
  \E Z_i^2
  \le
  (\E Z_i)^{(r-2)/(r-1)}(\E Z_i^r)^{1/(r-1)}.
\]
Hölder's inequality gives
\[
  \sum_i \E Z_i^2
  \le
  A^{(r-2)/(r-1)}B^{1/(r-1)}.
\]
Raising to \(r/2\) and using Young's inequality yields
\[
  \left(\sum_i \E Z_i^2\right)^{r/2}
  \le
  A^r+B.
\]
Combining this with
\(\norm{\sum_i Z_i}_r\le A+\norm{\sum_i(Z_i-\E Z_i)}_r\) proves the lemma.
\end{proof}

\begin{lemma}[Upper thinning]
\label{lem:upper-thin}
Let \(0\le a\le1\).  Let \(f_1,\ldots,f_n\in L^r(S)_+\) be deterministic, and
let \(R_1,\ldots,R_n\) be independent Bernoulli variables with
\(\Pp(R_i=1)\le a\).  Then
\[
  \E_R\PhiF\left(\sum_i R_i f_i\right)
  \le
  K_r\, a^{1/(2r)}
  \PhiF\left(\sum_i f_i\right).
\]
\end{lemma}

\begin{proof}
Let \(\rho_i=\Pp(R_i=1)\).  By Lemma \ref{lem:positive-rosenthal}, applied
pointwise in \(s\in S\),
\[
  \E_R\left(\sum_i R_i f_i(s)\right)^r
  \le
  C_r\left[
    \left(\sum_i \rho_i f_i(s)\right)^r
    +
    \sum_i \rho_i f_i(s)^r
  \right].
\]
Since \(\rho_i\le a\) and \(f_i\ge0\),
\[
  \left(\sum_i \rho_i f_i\right)^r\le a^r\left(\sum_i f_i\right)^r,
  \qquad
  \sum_i \rho_i f_i^r\le a\sum_i f_i^r\le a\left(\sum_i f_i\right)^r.
\]
Integrating and using \(a^r\le a\) gives
\[
  \E_R\norm{\sum_i R_i f_i}_{L^r}^r
  \le
  C'_r a\norm{\sum_i f_i}_{L^r}^r.
\]
Jensen's inequality for the concave map \(t\mapsto t^{\alpha/r}\) yields the
claim.
\end{proof}

\begin{lemma}[Lower thinning above one half]
\label{lem:lower-thin}
Let \(f_1,\ldots,f_n\in L^r(S)_+\) be deterministic.  Let
\(R_1,\ldots,R_n\) be independent Bernoulli variables with
\(\Pp(R_i=1)=c\), where \(1/2\le c\le1\).  Then
\[
  \PhiF\left(\sum_i f_i\right)
  \le
  2\,\E_R\PhiF\left(\sum_i R_i f_i\right).
\]
\end{lemma}

\begin{proof}
By subadditivity,
\[
  \PhiF\left(\sum_i f_i\right)
  \le
  \PhiF\left(\sum_i R_i f_i\right)
  +
  \PhiF\left(\sum_i (1-R_i)f_i\right).
\]
Taking expectations, the second term has the same law as a Bernoulli
\((1-c)\) thinning.  Since \(1-c\le c\), this thinning can be coupled as a
sub-thinning of a Bernoulli \(c\) thinning.  Monotonicity of \(\PhiF\) gives
\[
  \E\PhiF\left(\sum_i (1-R_i)f_i\right)
  \le
  \E\PhiF\left(\sum_i R_i f_i\right),
\]
and the result follows.
\end{proof}

\section*{Poissonization}

Let \(V_1,\ldots,V_n\) be independent \(L^r(S)_+\)-valued variables.  Write
\[
  p_i=\Pp(V_i\ne0).
\]
When \(p_i>0\), write \(Z_i\) for a random variable with the conditional law
of \(V_i\) given \(V_i\ne0\); when \(p_i=0\), omit that index.  Let
\[
  \Pi_V=\sum_i\sum_{k=1}^{N_i} Z_{i,k},
\]
where \(N_i\) are independent Poisson variables with means \(p_i\), and
\((Z_{i,k})\) are independent copies of \(Z_i\).  Equivalently, \(\Pi_V\) is
the compound-Poisson positive sum whose intensity measure is
\[
  \nu_V=\sum_i\mathcal L(V_i)|_{L^r_+\setminus\{0\}}.
\]

\begin{proposition}[Positive \(L^r\) Poissonization]
\label{prop:poisson}
For every \(r>1\), there is \(A_r<\infty\) such that
\[
  A_r^{-1}\,
  \E\PhiF\left(\sum_i V_i\right)
  \le
  \E\PhiF(\Pi_V)
  \le
  A_r\,
  \E\PhiF\left(\sum_i V_i\right).
\]
\end{proposition}

\begin{proof}
First prove the lower bound.  Let \(c=1-e^{-1}>1/2\).  Since
\[
  \Pp(N_i\ge1)=1-e^{-p_i}\ge c\,p_i
  \qquad(0\le p_i\le1),
\]
the first nonzero Poisson layer stochastically dominates a \(c\)-thinning of
the original family.  More explicitly, if \(V_i=B_iZ_i\) with
\(B_i\) Bernoulli \(p_i\), and if \(R_i\) are independent Bernoulli \(c\)
variables, then
\[
  \sum_i B_iR_iZ_i
\]
can be coupled below \(\sum_i \mathbf 1_{\{N_i\ge1\}}Z_{i,1}\), which is
itself dominated by \(\Pi_V\) in the positive cone.  Conditional on
\((B_i,Z_i)\), Lemma \ref{lem:lower-thin} gives
\[
  \PhiF\left(\sum_i B_iZ_i\right)
  \le
  2\,\E_R\PhiF\left(\sum_i B_iR_iZ_i\right).
\]
Taking expectations and using monotonicity gives
\[
  \E\PhiF\left(\sum_iV_i\right)
  \le
  2\,\E\PhiF(\Pi_V).
\]

For the upper bound, decompose the Poisson sum into layers:
\[
  \Pi_V
  =
  \sum_{k\ge1} L_k,
  \qquad
  L_k=\sum_i \mathbf 1_{\{N_i\ge k\}}Z_{i,k}.
\]
By subadditivity,
\[
  \E\PhiF(\Pi_V)\le \sum_{k\ge1}\E\PhiF(L_k).
\]
For \(0\le p\le1\),
\[
  \Pp(\operatorname{Pois}(p)\ge k)
  \le \frac{p^k}{k!}\le \frac{p}{k!}.
\]
The first inequality follows from
\(\mathbf 1_{\{N\ge k\}}\le \binom{N}{k}\) and
\(\E\binom{N}{k}=p^k/k!\) for \(N\sim\operatorname{Pois}(p)\).
Thus \(L_k\) has the law of a thinning of the original nonzero family with
retention probabilities at most \(a_k=1/k!\).  Applying Lemma
\ref{lem:upper-thin} conditionally to the original family gives
\[
  \E\PhiF(L_k)
  \le
  K_r\,(1/k!)^{1/(2r)}
  \E\PhiF\left(\sum_iV_i\right).
\]
Since \(\sum_{k\ge1}(1/k!)^{1/(2r)}<\infty\), summing over \(k\) proves the
upper bound.
\end{proof}

\begin{theorem}[Positive aggregate comparison for \(q>2\)]
\label{thm:positive-qgt2}
Let \(q>2\), \(r=q/2\).  Let \(V_i,W_j\) be finite independent families of
positive \(L^r(S)\)-valued random variables.  If
\[
  \sum_i\mathcal L(V_i)|_{L^r_+\setminus\{0\}}
  \le
  \sum_j\mathcal L(W_j)|_{L^r_+\setminus\{0\}},
\]
then
\[
  \E\norm{\sum_iV_i}_{L^r}^{1/2}
  \le
  C_q
  \E\norm{\sum_jW_j}_{L^r}^{1/2}.
\]
\end{theorem}

\begin{proof}
Let \(\nu_V\le\nu_W\) be the two aggregate occurrence measures.  Construct
compound-Poisson sums \(\Pi_V,\Pi_W\).  Since \(\nu_W=\nu_V+(\nu_W-\nu_V)\),
\(\Pi_W\) has the same law as \(\Pi_V+\Pi'\), where \(\Pi'\) is an independent
positive compound-Poisson sum.  Monotonicity of \(\PhiF\) gives
\[
  \E\PhiF(\Pi_V)\le \E\PhiF(\Pi_W).
\]
Apply Proposition \ref{prop:poisson} to both families.
\end{proof}

\section*{Proof of the Square-Function Theorem}

\begin{proof}[Proof of Theorem \ref{thm:square}]
For \(1<q\le2\), this is exactly the theorem proved in the companion packet
\[
\texttt{1907.11588\_Lq\_p1\_endpoint\_qle2\_square\_function}.
\]

Let \(q>2\) and set \(r=q/2\).  Given \(L^q(S)\)-valued variables \(x_i,y_j\),
put
\[
  V_i=|x_i|^2,\qquad W_j=|y_j|^2.
\]
The map \(x\mapsto |x|^2\) is continuous from \(L^q(S)\) to \(L^{q/2}(S)\),
so aggregate domination \(\mu_x\le\mu_y\) pushes forward to aggregate
domination of the positive \(L^r\)-valued square magnitudes.  Theorem
\ref{thm:positive-qgt2} gives
\[
  \E\norm{\sum_i |x_i|^2}_{L^r}^{1/2}
  \le
  C_q
  \E\norm{\sum_j |y_j|^2}_{L^r}^{1/2},
\]
which is precisely \(\mathsf{SF}_q(x)\le C_q\mathsf{SF}_q(y)\).
\end{proof}

\section*{Discrete and General-Local Consequences}

\begin{proposition}[Discrete \(p=1\) comparison]
\label{prop:discrete}
Let \(1<q<\infty\).  Let \((\xi_i)_{i=1}^n\) and \((\eta_j)_{j=1}^m\) be
independent mean-zero \(L^q(S)\)-valued random variables.  If
\[
   \sum_i \Pp(\xi_i\in B)
   \le
   \sum_j \Pp(\eta_j\in B)
   \qquad(B\in\cB(\Lq),\;0\notin B),
\]
then
\[
   \E\max_{1\le k\le n}
      \left\|\sum_{i=1}^k \xi_i\right\|_{L^q(S)}
   \le
   C_q
   \E\max_{1\le k\le m}
      \left\|\sum_{j=1}^k \eta_j\right\|_{L^q(S)}.
\]
\end{proposition}

\begin{proof}
The aggregate hypothesis is \(\mu_\xi\le\mu_\eta\), so Theorem
\ref{thm:square} compares the two lattice square functions.  Yaroslavtsev's
Banach-function-space BDG theorem
\cite[Theorem 4.32 and Proposition 4.33]{Yaroslavtsev2018} says that for
every \(L^q(S)\)-valued martingale difference sequence \(d=(d_i)\),
\[
  \E\sup_N\left\|\sum_{i=1}^N d_i\right\|_{L^q}
  \simeq_q
  \E\left\|
      \left(\sum_i |d_i|^2\right)^{1/2}
    \right\|_{L^q}.
\]
Both independent mean-zero families are martingale difference sequences in
any fixed order.  Applying this equivalence to \((\xi_i)\) and \((\eta_j)\)
proves the maximal comparison.
\end{proof}

\begin{proof}[Proof of Theorem \ref{thm:main}]
The conditional endpoint packet in this run proves that Theorem
\ref{thm:square}, for a fixed \(q\), implies the full \(L^q(S)\)-valued
general-local characteristic-domination theorem at \(p=1\).  For completeness
we recall the reduction.

Yaroslavtsev's decoupled tangent construction and canonical decomposition
reduce the comparison to continuous parts and purely discontinuous tangent
jump martingales.  Continuous parts are controlled at \(p=1\) by the source's
continuous characteristic-domination estimate \cite[Proposition 12.6]{Yaroslavtsev2019}.
For the jump part, conditional on the local characteristics and after the
usual truncation of time and jump size, the decoupled tangent jumps are finite
independent mean-zero families.  Characteristic domination controls the total
spatial jump compensator,
\[
  \nu^N(\R_+\times B)\le \nu^M(\R_+\times B),
  \qquad B\in\cB(L^q(S)),\ 0\notin B,
\]
which is exactly the aggregate occurrence-measure domination in Proposition
\ref{prop:discrete}.

Apply Proposition \ref{prop:discrete} conditionally and pass through the
truncation and localization limits.  Combining the discontinuous estimate
with the continuous part and the \(p=1\) Banach-function-space
BDG/canonical-decomposition equivalence \cite[Section 4.3]{Yaroslavtsev2018}
gives the zero-initial-value case.  Nonzero initial values are handled by the
initial-value domination in Definition 12.9 and the triangle inequality.
\end{proof}

\section*{Verification and Novelty Notes}

The new argument is self-contained apart from scalar Rosenthal's inequality
and Yaroslavtsev's \(p=1\) Banach-function-space BDG theorem.  The
Montgomery-Smith--Pruss comparison theorem \cite{MontgomerySmithPruss} was
used as a guide for looking at iid and Poisson aggregate models, but the
Poissonization proposition above is proved directly for the positive
\(L^r\)-lattice functional needed here.

Cheap run indexes contained prior packets for \(1<p\) in \(L^q\), Hilbert all
\(p\), \(\ell_q\) at \(p=q\), the \(1<q\le2\) endpoint, a deterministic
\(q>2\) subcase, the broad \(c_0\) counterexample, and the conditional
endpoint reduction.  No packet proved the full \(L^q(S)\), \(p=1\),
\(1<q<\infty\) endpoint.

\section*{Human Review Recommendation}

Review Proposition \ref{prop:poisson} first.  The lower bound uses only the
fact that \(1-e^{-p}\ge(1-e^{-1})p\) and the lower thinning lemma.  The upper
bound is the only delicate step: decompose Poisson counts into layers and use
the upper thinning lemma with the summable factors \((1/k!)^{1/q}\).  Once
that proposition is accepted, aggregate domination follows immediately by
Poisson superposition.

\begin{thebibliography}{9}

\bibitem{Yaroslavtsev2019}
I. S. Yaroslavtsev.
\newblock Local characteristics and tangency of vector-valued martingales.
\newblock arXiv:1907.11588, 2019.

\bibitem{Yaroslavtsev2018}
I. S. Yaroslavtsev.
\newblock Burkholder--Davis--Gundy inequalities in UMD Banach spaces.
\newblock arXiv:1807.05573, 2018.

\bibitem{DirksenYaroslavtsev2017}
S. Dirksen and I. S. Yaroslavtsev.
\newblock \(L^q\)-valued Burkholder--Rosenthal inequalities.
\newblock arXiv:1707.00109, 2017.

\bibitem{Rosenthal1970}
H. P. Rosenthal.
\newblock On the subspaces of \(L_p\), \(p>2\), spanned by sequences of
  independent random variables.
\newblock Israel Journal of Mathematics 8 (1970), 273--303.

\bibitem{MontgomerySmithPruss}
S. J. Montgomery-Smith and A. R. Pruss.
\newblock A comparison inequality for sums of independent random variables.
\newblock arXiv:math/9811124, 1998.

\end{thebibliography}

\end{document}
